\declareWideFig{primitives}{A unified framework for KV management}{(a) \textit{KV Admission (Ours)} predictively filters redundant tokens pre-write. (b) \textit{KV Selection} selectively attends to relevant tokens at read-time. (c) \textit{KV Eviction} retrospectively prunes obsolete tokens post-write.}

\declareFig{breakdown}{Attention bottleneck in long-context LLM inference}{Measured on Llama-3.1-8B on an H200 GPU.}

\declareFig{distribution}{Heterogeneity of token utility}{}

\declareFig{fragmentation}{Memory fragmentation with ragged KV states}{}

\declareFig{overview}{Overview of Write-Gated KV}{}

\declareWideFig{management}{Dual-cache paged memory management}{(a--c) To handle ragged cache growth across heads, we decouple the logical view into a fixed-size \textit{local cache} and a growing \textit{global cache}, mapped to a unified physical \textit{KV pool} via per-head \textit{page tables} to prevent memory fragmentation. (d) The cache update logic during decoding. We employ \textit{lazy promotion} where the token exiting the local window is inspected (and possibly promoted) upon replacement.}

\declareWideFig{benchmark}{Long-context HELMET task accuracy on Llama-3.1-8B}{The x-axis represents the normalized KV cache size. WG-KV consistently outperforms baselines under tight admission budgets.}

\declareFig{perf}{End-to-end latency and memory usage on Llama-3.1-8B}{}

\declareRoundedFig{interpret_short}{5pt}{Visualization of admission decisions}{Highlighted tokens are admitted to the global cache.}

\declareWideFig{benchmark_qwen}{Long-context HELMET task accuracy on Qwen3-4B-2507}{The trends align with the results on Llama-3.1-8B, showing WG-KV's generalizability across model architectures.}

\declareFig{perf_qwen}{End-to-end latency and memory usage on Qwen3-4B-2507}{}

\declareWideFig{quest}{Synergy between WG-KV and Quest}{By pre-filtering redundant tokens, WG-KV effectively (a) reduces KV memory usage. This compacted cache shrinks the candidate pool for selection, leading to (b) lower decode-time memory access and (c) reduced attention latency. Furthermore, integrating WG-KV (d) maintains nearly identical task accuracy compared to standalone Quest.}

\declareWideFig{integration_new}{Long-context HELMET task accuracy under varying selection budgets}{The results demonstrate that integrating WG-KV with Quest consistently matches the accuracy of standalone Quest, highlighting the composable nature of pre-write Admission and read-time Selection.}

\declareFig{evict}{Conceptual illustration of the synergy between KV Admission and Eviction}{}

\declareFig{snapkv_short}{Synergy between WG-KV and SnapKV on AIME25 under strict memory bounds}{}

\declareFig{snapkv}{Admission-Eviction synergy on AIME25 using Qwen3-4B-2507}{``Off'' denotes the baseline where WG-KV is disabled, and the reported metrics (KV Cache Size and \# Eviction Triggers) are averaged across all samples.}

\declareRoundedFig{evict_interpret}{5pt}{Visualization of the token lifecycle during an AIME25 reasoning trace using Qwen3-4B-2507 on Layer 25, Head 3}{}

\declareFig{threshold}{Impact of $\lambda$ and $\tau$ on sparsity-loss tradeoff}{}

\declareFig{threshold_cmp}{Impact of local cache on sparsity-loss tradeoff}{Removing local cache leads to a sharp surge in distillation loss.}

\declareWideFig{gate}{Visualization of input-dependent admission patterns}{The heatmaps show the normalized KV cache size across each head for Llama-3.1-8B with WG-KV. We compare the \textit{(a) Code Summarization} task, which exhibits a relatively uniform distribution, against the \textit{(b) HTML to TSV} task, which shows a sparser pattern. This contrast demonstrates that WG-KV learns an input-dependent policy based on the semantic structure of the task.}

\declareRoundedWideFig{l0h7}{10pt}{Full visualization of admission decisions for Layer 1, Head 8}{}

\declareRoundedWideFig{l4h7}{10pt}{Full visualization of admission decisions for Layer 5, Head 8}{}

\declareRoundedWideFig{l24h0}{10pt}{Full visualization of admission decisions for Layer 25, Head 1}{}

\declareRoundedWideFig{l27h2}{10pt}{Full visualization of admission decisions for Layer 28, Head 3}{}

\declareRoundedWideFig{l28h7}{10pt}{Full visualization of admission decisions for Layer 29, Head 8}{}

\declareRoundedWideFig{l13h1}{10pt}{Full visualization of admission decisions for Layer 14, Head 2}{}

\declareRoundedWideFig{l18h6}{10pt}{Full visualization of admission decisions for Layer 19, Head 7}{}

\declareRoundedWideFig{interpret_summary}{10pt}{Aggregated admission decisions across all attention heads}{Color intensity denotes the proportion of heads admitting each token. High cross-head consensus is observed on structural control tokens, key entities, and semantically dense words, whereas function words and pronouns exhibit lower admission ratios.}

\newcommand{\tabReplay}[2]{
  \begin{table}[#1]
    \centering
    \def\captext{Average last-layer hidden states MSE under different admission policies}
    \caption[\captext]{\captext.}
    \label{tab:replay}
    \resizebox{#2\linewidth}{!}{
      \begin{tabular}{ccc}
        \toprule
        \textbf{Own Policy} & \textbf{Replay Policy} & \textbf{Random Policy} \\ 
        \midrule
        0.01480 & 0.10297 & 0.10271 \\ 
        \bottomrule
      \end{tabular}
    }
  \end{table}
}